% Source: originals/shrock-phy541/homework/phy541_hw1.pdf, pages 1-2 (photographed handout dated 2/11/2026).
% Placed after the 1D Ising model; it uses Lectures 1-7.
\hwset{1}

\noindent\emph{Dated 2/11/2026.} Please submit your homework papers at the beginning of class on Mon.\ Feb.\ 23. It is
not necessary to submit typed answers, but if you write your answers by hand, please make your handwriting legible.
This assignment has a total of 120 points.

\begin{hwproblem}{1}{1}
(20 pts) Consider the air outside on a winter night at a temperature of $-13^\circ$ Celsius and pressure of one
atmosphere. (a) (10 pts) Calculate the mean distance between the gas molecules. (b) (5 pts) How does this distance
compare with the approximate sizes of the N$_2$ and O$_2$ molecules in the air? (c) (5 pts) Explain how this is
related to the accuracy of the ideal gas law at this temperature and pressure.
\end{hwproblem}

\begin{hwproblem}{1}{2}
(20 pts) Consider a spin model with spins $\mathbf S_i$ at sites $i$ on (the thermodynamic limit of) a lattice
$\Lambda$ and a spin-spin interaction given by the Hamiltonian $\Ham = -J\sum_{e_{i\ell}}\mathbf S_i\cdot\mathbf S_\ell$
with $J>0$, i.e., a ferromagnetic spin-spin coupling. Here, $e_{i\ell}$ denotes the bond between adjacent (i.e.,
nearest-neighbor) sites $i$ and $\ell$ on $\Lambda$. The system is in thermal equilibrium at a temperature $T$. Let
$\beta = 1/(\kB T)$, $K = \beta J$, and $u = e^{-4K}$. Assume that the spontaneous magnetization (per site) in this model
is given by
\begin{equation}
  M = \frac{(1+u)^{1/4}(1-6u+u^2)^{1/7}}{(1-u)^{1/2}}
\end{equation}
for $T<T_c$ (and zero for $T\ge T_c$), where $T_c$ is the critical temperature. Calculate (a) (10 pts) the critical
temperature $T_c$; (b) (10 pts) the critical exponent $\beta$.
\end{hwproblem}

\begin{hwproblem}{1}{3}
(20 pts) Assume that the initial susceptibility $\chi_0$ of a given magnetic system has the following form, for
temperatures $T>T_c$, as $T-T_c\to 0^+$:
\begin{equation}
  \chi_0 = \frac{b}{\left[(T/T_c)^8-1\right]^{3/2}}
\end{equation}
where $b$ is a (positive) real number. Calculate the exponent $\gamma$, which gives the leading singularity in
$\chi_0$ as $T\to T_c^+$ (20 pts).
\end{hwproblem}

\begin{hwproblem}{1}{4}
(30 pts) Consider the spin-$1/2$ Ising model in 1D in an external (uniform) magnetic field, $H$, described by the
Hamiltonian
\begin{equation}
  \Ham = -J\sum_{e_{ij}}\sigma_i\sigma_j - H\sum_i\sigma_i
\end{equation}
and partition function
\begin{equation}
  Z = \sum_{\{\sigma_i\}}e^{-\beta\Ham}
\end{equation}
where $\sigma_i = \pm1$ at each site $i$ and $e_{ij}$ denotes the edge (bond) connecting adjacent sites $i$ and $j$.
As in class, define $K = \beta J$ and $h = \beta H$, where $\beta = 1/(\kB T)$. In class we calculated the induced
magnetization per site as $M(H)$ as
\begin{equation}
  M = \frac{\sinh h}{\left[\sinh^2h+e^{-4K}\right]^{1/2}}.
\end{equation}
Consider the case of $J<0$, i.e., antiferromagnetic spin-spin coupling. Calculate $M$ at $T = 0$. Hint: This is a
case of competing interactions.
\end{hwproblem}

\begin{hwproblem}{1}{5}
(20 pts) The spin-$s$ Ising model is defined as a spin model with classical spin variables $\sigma_i$ that can take on
values in the set $\{-s,-s+1,\ldots,s\}$. The Hamiltonian and partition function are given by
\begin{gather}
  \Ham = -J\sum_{\{e_{ij}\}}\sigma_i\sigma_j - H\sum_i\sigma_i\\
  Z = \sum_{\{\sigma_i\}}e^{-\beta\Ham} = \sum_{\{\sigma_i\}}e^{K\sum_{\{e_{ij}\}}\sigma_i\sigma_j+h\sum_i\sigma_i}
\end{gather}
where again $e_{ij}$ denotes the bond (edge) connecting two nearest-neighbor sites, and $K = \beta J$, and $h = \beta H$.
This problem deals with the spin-1 Ising model, where spins at each site take on values in the set $\{-1,0,1\}$. Let us
consider this model on a 1D lattice with periodic boundary conditions. (a) (10 pts) Calculate the transfer matrix for
this model on the given lattice. (b) (10 pts) Calculate the Gibbs free energy per site in the case of zero external
field, $H = 0$, in the thermodynamic limit.
\end{hwproblem}

\begin{hwproblem}{1}{6}
(20 pts) Determine the planar dual of the $3\cdot 6\cdot 3\cdot 6$ lattice. Is this an Archimedean lattice? Prove
your answer.
\end{hwproblem}
